\chapter{Differentiating Through the Forward Solver}\label{ch:adjoint}

The previous chapter postulated that gradient information is could be desirable
both for initialising the Markov chain and for
constructing more efficient MCMC proposals. Computing the gradient of the
log-likelihood~\cref{eq:log-likelihood} with respect to the
parameters~$\phi$ requires differentiating through the forward
map~$\mathcal{P}$, which itself involves solving a sequence of tridiagonal
linear systems. In this chapter we describe how these derivatives are
obtained efficiently.

\section{Discretise Then Optimise}\label{sec:disc-then-opt}

There are broadly two strategies for differentiating through a numerical
solver. In the \emph{optimise then discretise} approach one derives the
continuous adjoint equations analytically and then discretises them; in the
\emph{discretise then optimise} approach one first discretises the forward
problem and then differentiates the discrete computation. Generally,
the discretise then optimise is preferable as it computes the exact
gradient of the quantity that is actually evaluated in code, avoiding
the consistency errors that can arise
when the adjoint PDE is discretised independently of the forward PDE
\cite{kidger2021neural}. We
adopt this approach throughout.

Concretely, our forward map~$\mathcal{P}$ is implemented as a
sequence of algebraic operations --- assembling tridiagonal systems and
solving them via the Thomas algorithm at each timestep, then extracting
the surface flux. The gradient we seek is therefore the derivative of a
composition of these discrete operations with respect to the
parameters~$\phi$.

\section{Reverse Mode Automatic Differentiation}\label{sec:reverse-ad}

The log-likelihood is a scalar-valued function of the parameter
vector. In this setting  reverse mode
automatic differentiation (AD) is the
natural choice, as it computes the full gradient in a single backward
pass regardless of the number of parameters.

Reverse mode AD decomposes the gradient computation into a sequence of
\emph{vector--Jacobian products} (VJPs). If the forward computation is
written as a composition $f = f_L \circ f_{L-1} \circ \cdots \circ f_1$,
then the chain rule gives
\begin{equation}
  \frac{\partial f}{\partial \phi}
  = \frac{\partial f_L}{\partial f_{L-1}}
  \cdot \frac{\partial f_{L-1}}{\partial f_{L-2}}
  \cdots
  \frac{\partial f_1}{\partial \phi}.
  \label{eq:chain-rule}
\end{equation}
In reverse mode, this product is evaluated from left to right: at each
stage a row vector (the adjoint, or cotangent) is multiplied by the
Jacobian of the corresponding primitive operation. The intermediate
Jacobians are never formed explicitly; only the VJP of each primitive is
required.

For the forward map~$\mathcal{P}$ the dominant primitives are the
tridiagonal solves --- one per timestep. The question therefore reduces
to: how should the VJP of a tridiagonal solve be implemented?

\section{Differentiating Through Linear Solves}\label{sec:diff-linear}

Consider a single linear system $A x = b$, where $A$ depends on the
parameters~$\phi$ and
$b$ depends on the solution at the previous timestep. A na\"ive approach
is to record every intermediate value of the Thomas algorithm on the
forward pass and replay them in reverse during the backward pass. This is
what a generic AD framework will do if the Thomas algorithm is
implemented as a sequence of elementary operations: the entire
computational graph of the factorisation and back-substitution is stored
on a tape.

For a single tridiagonal solve on $n$ unknowns the memory overhead of
taping is modest. However, the forward map involves $m$ consecutive
solves (one per timestep), and the tape must be retained across all of
them because the backward pass traverses the full time history in
reverse. The total memory cost therefore scales as $O(mn)$, which becomes
prohibitive for fine spatial and temporal grids. Worse, the overhead of
maintaining and traversing the tape degrades runtime performance even when
memory is not the binding constraint.

\section{Adjoint of the Tridiagonal Solve}\label{sec:adjoint-solve}

A more efficient approach exploits the structure of the linear solve
directly. Suppose the forward pass solves $Ax = b$ and the backward pass
receives the adjoint (cotangent) $\bar{x} = \partial \ell / \partial x$.
We require the adjoint with respect to~$b$:
\begin{equation}
  \bar{b} = \frac{\partial \ell}{\partial b}
  = \bar{x}^T \frac{\partial x}{\partial b}
  = \bar{x}^T A^{-1}.
  \label{eq:adjoint-b}
\end{equation}
This is equivalent to solving the \emph{adjoint system}
\begin{equation}
  A^T \lambda = \bar{x},
  \label{eq:adjoint-system}
\end{equation}
and setting $\bar{b} = \lambda^T$. Similarly, the adjoint with respect to
the matrix entries of~$A$ is
\begin{equation}
  \bar{A} = -\lambda\, x^T,
  \label{eq:adjoint-A}
\end{equation}
which follows from differentiating the identity $A x = b$ and
substituting the solution of the adjoint system.

The key observation is that $A^T$ is also tridiagonal (with the sub- and
super-diagonals transposed), so the adjoint system~\cref{eq:adjoint-system}
can be solved by the Thomas algorithm at the same $O(n)$ cost as the
forward solve. Crucially, this approach requires only the matrix~$A$ and
the forward solution~$x$ to be stored --- not the full tape of
intermediate values from the factorisation. The memory cost per timestep
is therefore $O(n)$ rather than the $O(n)$ per operation of na\"ive
taping, and the total cost across all $m$ timesteps is reduced from
taping the entire forward history to storing only the tridiagonal
coefficients and solutions at each level.

In practice we implement this by defining a custom VJP rule for the
tridiagonal solve in \emph{JAX} \cite{jax2018github}. On the forward
pass the Thomas algorithm solves
$Ax = b$ and caches $A$ and~$x$. On the backward pass the custom rule
receives~$\bar{x}$, solves $A^T \lambda = \bar{x}$ using the banded
transpose, and returns $\bar{b} = \lambda^T$ and
$\bar{A} = -\lambda\, x^T$. No taping of the Thomas algorithm's internal
operations is required.

\section{Extension to Pentadiagonal Systems}\label{sec:pentadiagonal}

The benefits of the custom adjoint become even more pronounced for
higher-order discretisations that produce pentadiagonal (or wider banded)
systems. A pentadiagonal factorisation involves substantially more
intermediate values than a tridiagonal one, and taping the full forward
computation through the banded solver would incur a correspondingly
larger memory and runtime penalty. The adjoint approach sidesteps this
entirely: the transpose of a pentadiagonal matrix is again
pentadiagonal, so the adjoint system is solved by a banded solver at the
same $O(n)$ cost per timestep. The custom VJP rule generalises directly
--- one simply solves with the transposed band structure --- making this
a crucial step in enabling the use of higher-order spatial
discretisations within the differentiable forward map.

\section{Initialisation Results}
